\chapter{Einleitung}
Online-Produktrezensionen sind eine zentrale Informationsquelle im E-Commerce und 
beeinflussen maßgeblich Kaufentscheidungen. Sie enthalten wertvolle Rückmeldungen 
zu Produktqualität und Nutzererfahrungen und bilden eine umfangreiche Datenbasis 
für Unternehmen. Um diese Informationen systematisch auszuwerten, wird oft eine
Sentimentanalyse verwendet. Dieses Verfahren extrahiert Meinungen und Emotionen
aus Texten und klassifiert sie hinsichtlich ihrer Polarität.
Die Bewertung erfolgt häufig über Klassen oder numerische 
Skalen [4].

\section{Problemstellung}
Online-Rezensionen liegen in großer Anzahl vor und sind aufgrund ihrer 
Sprachvielfalt und subjektiven Ausdrucksweise nur schwer manuell auswertbar. Eine 
systematische Analyse ist zeitaufwendig und nicht sinnvoll skalierbar.
Automatische Sentimentanalyse-Systeme müssen mehrdeutige Formulierungen, implizite 
Bewertungen und unterschiedliche sprachliche Ausdrücke korrekt verarbeiten. Dabei 
sind eine eindeutige Repräsentation von Meinungen, die Zusammenführung semantisch 
ähnlicher Aussagen sowie die Ableitung nicht explizit genannter Informationen 
erforderlich [3]. Zudem sind fein abgestufte Bewertungen, wie Sternebewertungen, 
oft aussagekräftiger als binäre Klassifikationen [6].

\section{Zielsetzung}
Ziel dieser Arbeit ist die Entwicklung eines automatischen Textklassifikationssystems 
zur Sentimentanalyse von Online-Produktrezensionen. Das System soll in der Lage sein, 
Textbewertungen einer numerischen Skala von 1 bis 5 Sternen zuzuordnen und damit eine 
Mehrklassen-Klassifikation durchzuführen. Der Fokus liegt auf der Nutzung 
maschineller Lernverfahren, da diese im Vergleich zu regelbasierten Ansätzen besser 
mit sprachlicher Vielfalt, Tippfehlern und unbekannten Ausdrücken umgehen können und 
übergeordnete Sinnzusammenhänge erkennen [6].
Zu diesem Zweck werden zwei etablierte Klassifikationsmodelle eingesetzt: 
logistische Regression und Multinomial Naive Bayes. Ziel ist es, die 
Leistungsfähigkeit dieser beiden Ansätze zu vergleichen und ihre 
Eignung für die Vorhersage fein abgestufter Bewertungen zu bewerten.

\section{Vorgehensweise}
Zu Beginn der Arbeit erfolgt eine explorative Datenanalyse, um einen 
Überblick über den Datensatz, dessen Struktur sowie zentrale Eigenschaften zu 
gewinnen. Darauf aufbauend wird die Datenvorverarbeitung beschrieben, die unter 
anderem Schritte zur Textbereinigung und Merkmalsextraktion umfasst.
Anschließend werden die eingesetzten Klassifikationsverfahren vorgestellt und 
methodisch eingeordnet. Im Rahmen des Modelltrainings werden eine logistische 
Regression sowie ein Multinomial-Naive-Bayes-Klassifikator implementiert und 
trainiert.
Die Leistungsfähigkeit der Modelle wird mithilfe geeigneter Evaluationsmetriken 
bewertet und miteinander verglichen. Abschließend werden die Ergebnisse 
diskutiert und im Hinblick auf die Stärken und Schwächen der jeweiligen Ansätze 
analysiert.
